% GENERATED FILE. Edit canonical lessons, then run npm run generate:books.
% canonical-source-hash: fnv1a64:b7674061b8399ef2
% canonical-lessons: KA-C210-konege, KA-C210-samanyavagi, KA-C210-taksana, KA-C210-biduvu, KA-C210-yugadi

\chapter{Finally, Usually, Immediately, Free time, Ugadi}
\label{ch:ka-finally-usually-immediately-free-time-ugadi}
\hlchaptermodality{210}

\begin{chapteropening}
\textbf{By the end of this chapter:} \emph{I can say finally, usually, immediately, free time and Ugadi.}

Complete the last lesson of chapter 210: I can say finally, usually, immediately, free time and Ugadi.
\end{chapteropening}

\section[\texorpdfstring{\kn{ಕೊನೆಗೆ} (konege)}{ಕೊನೆಗೆ (konege)}]{\kn{ಕೊನೆಗೆ} (konege) — finally, in the end}

\label{lesson:KA-C210-konege}



\begin{quote}
Before the new one: say the Kannada for yoga, then the Kannada for meditation.
\end{quote}

\subsection*{You'll want to know: \kn{ಕೊನೆಗೆ}}
\textbf{\kn{ಕೊನೆಗೆ}} — \emph{konege} — \textquotedblleft{}finally, in the end\textquotedblright{}.

\begin{grammarlens}[title={What you've built}]
One more word for this chapter.
\end{grammarlens}

\subsection*{Your turn}
\begin{itemize}
\raggedright
  \item \emph{Say it:} \emph{konege}
  \item \emph{Say it:} \emph{konege}, once more
  \item \emph{Say it:} say \emph{konege}
  \item \emph{Recall:} say the Kannada for yoga, then the Kannada for meditation, then say \emph{konege} again
\end{itemize}

\begin{tcolorbox}[breakable,colback=teal!4,colframe=teal!35!black,title={Before you move on}]
What does \kn{ಕೊನೆಗೆ} mean? (\textquotedblleft{}Finally, in the end\textquotedblright{}.) Say it once more.
\end{tcolorbox}
\section[\texorpdfstring{\kn{ಸಾಮಾನ್ಯವಾಗಿ} (sāmānyavāgi)}{ಸಾಮಾನ್ಯವಾಗಿ (sāmānyavāgi)}]{\kn{ಸಾಮಾನ್ಯವಾಗಿ} (sāmānyavāgi) — usually}

\label{lesson:KA-C210-samanyavagi}



\begin{quote}
Before the new one: say the Kannada for meditation, then the Kannada for finally.
\end{quote}

\subsection*{You'll want to know: \kn{ಸಾಮಾನ್ಯವಾಗಿ}}
\textbf{\kn{ಸಾಮಾನ್ಯವಾಗಿ}} — \emph{sāmānyavāgi} — \textquotedblleft{}usually\textquotedblright{}.

\begin{grammarlens}[title={What you've built}]
One more word for this chapter.
\end{grammarlens}

\subsection*{Your turn}
\begin{itemize}
\raggedright
  \item \emph{Say it:} \emph{sāmānyavāgi}
  \item \emph{Say it:} \emph{sāmānyavāgi}, once more
  \item \emph{Say it:} say \emph{sāmānyavāgi}
  \item \emph{Recall:} say the Kannada for meditation, then the Kannada for finally, then say \emph{sāmānyavāgi} again
\end{itemize}

\begin{tcolorbox}[breakable,colback=teal!4,colframe=teal!35!black,title={Before you move on}]
What does \kn{ಸಾಮಾನ್ಯವಾಗಿ} mean? (\textquotedblleft{}Usually\textquotedblright{}.) Say it once more.
\end{tcolorbox}
\section[\texorpdfstring{\kn{ತಕ್ಷಣ} (takṣaṇa)}{ತಕ್ಷಣ (takṣaṇa)}]{\kn{ತಕ್ಷಣ} (takṣaṇa) — immediately}

\label{lesson:KA-C210-taksana}



\begin{quote}
Before the new one: say the Kannada for finally, then the Kannada for usually.
\end{quote}

\subsection*{You'll want to know: \kn{ತಕ್ಷಣ}}
\textbf{\kn{ತಕ್ಷಣ}} — \emph{takṣaṇa} — \textquotedblleft{}immediately\textquotedblright{}.

\begin{grammarlens}[title={What you've built}]
One more word for this chapter.
\end{grammarlens}

\subsection*{Your turn}
\begin{itemize}
\raggedright
  \item \emph{Say it:} \emph{takṣaṇa}
  \item \emph{Say it:} \emph{takṣaṇa}, once more
  \item \emph{Say it:} say \emph{takṣaṇa}
  \item \emph{Recall:} say the Kannada for finally, then the Kannada for usually, then say \emph{takṣaṇa} again
\end{itemize}

\begin{tcolorbox}[breakable,colback=teal!4,colframe=teal!35!black,title={Before you move on}]
What does \kn{ತಕ್ಷಣ} mean? (\textquotedblleft{}Immediately\textquotedblright{}.) Say it once more.
\end{tcolorbox}
\section[\texorpdfstring{\kn{ಬಿಡುವು} (biḍuvu)}{ಬಿಡುವು (biḍuvu)}]{\kn{ಬಿಡುವು} (biḍuvu) — free time, leisure}

\label{lesson:KA-C210-biduvu}



\begin{quote}
Before the new one: say the Kannada for usually, then the Kannada for immediately.
\end{quote}

\subsection*{You'll want to know: \kn{ಬಿಡುವು}}
\textbf{\kn{ಬಿಡುವು}} — \emph{biḍuvu} — \textquotedblleft{}free time, leisure\textquotedblright{}.

\begin{grammarlens}[title={What you've built}]
One more word for this chapter.
\end{grammarlens}

\subsection*{Your turn}
\begin{itemize}
\raggedright
  \item \emph{Say it:} \emph{biḍuvu}
  \item \emph{Say it:} \emph{biḍuvu}, once more
  \item \emph{Say it:} say \emph{biḍuvu}
  \item \emph{Recall:} say the Kannada for usually, then the Kannada for immediately, then say \emph{biḍuvu} again
\end{itemize}

\begin{tcolorbox}[breakable,colback=teal!4,colframe=teal!35!black,title={Before you move on}]
What does \kn{ಬಿಡುವು} mean? (\textquotedblleft{}Free time, leisure\textquotedblright{}.) Say it once more.
\end{tcolorbox}
\section[\texorpdfstring{\kn{ಯುಗಾದಿ} (yugādi)}{ಯುಗಾದಿ (yugādi)}]{\kn{ಯುಗಾದಿ} (yugādi) — Ugadi, the Kannada New Year}

\label{lesson:KA-C210-yugadi}



\begin{quote}
Before the new one: say the Kannada for immediately, then the Kannada for free time.
\end{quote}

\subsection*{You'll want to know: \kn{ಯುಗಾದಿ}}
\textbf{\kn{ಯುಗಾದಿ}} — \emph{yugādi} — \textquotedblleft{}Ugadi, the Kannada New Year\textquotedblright{}.

\begin{grammarlens}[title={What you've built}]
One more word for this chapter.
\end{grammarlens}

\subsection*{Your turn}
\begin{itemize}
\raggedright
  \item \emph{Say it:} \emph{yugādi}
  \item \emph{Say it:} \emph{yugādi}, once more
  \item \emph{Say it:} say \emph{yugādi}
  \item \emph{Recall:} say the Kannada for immediately, then the Kannada for free time, then say \emph{yugādi} again
\end{itemize}

\begin{tcolorbox}[breakable,colback=teal!4,colframe=teal!35!black,title={Before you move on}]
What does \kn{ಯುಗಾದಿ} mean? (\textquotedblleft{}Ugadi, the Kannada New Year\textquotedblright{}.) Say it once more.
\end{tcolorbox}
